% Generated mechanically from frozen input p12/ja/fdl.tex.
% Do not edit this generated file; edit the bound locale source instead.

% OK, start here.
%
%---------------------------------------------------------------------
%\tableofcontents

\setcounter{chapter}{115}
\chapter{GNU 自由文書ライセンス（GNU Free Documentation License）}

\phantomsection
\label{fdl-section-phantom}

\begin{center}

      バージョン 1.2、2002年11月


Copyright \copyright 2000, 2001, 2002 Free Software Foundation, Inc.

\bigskip

    51 Franklin St, Fifth Floor, Boston, MA  02110-1301  USA

\bigskip

Everyone is permitted to copy and distribute verbatim copies
of this license document, but changing it is not allowed.

本ライセンス文書を一字一句そのままに複製し、頒布することは誰にでも許可されるが、
変更することは認められない。
\end{center}

This is an unofficial translation of the GNU Free Documentation License into
Japanese. It was not published by the Free Software Foundation, and does not
legally state the distribution terms for documentation that uses the GNU
FDL---only the original English text of the GNU FDL does that. However, we hope
that this translation will help Japanese speakers understand the GNU FDL
better.

以下は GNU Free Documentation License の非公式な日本語参考訳である。これは Free
Software Foundation が発表したものではなく、GNU FDL を適用する文書の頒布条件を法的に
有効な形で定めるものでもない。その条件を定めるのは GNU FDL の英語原文だけである。
本参考訳は、八田真行による2005年11月9日付の非公式日本語訳を照合資料として英語原文に
合わせて再構成したものであり、日本語話者が GNU FDL を理解するための便宜として提供する。
本参考訳は英語原文と共に提供し、両者に不一致がある場合は英語原文を優先する。

You may publish this translation, modified or unmodified, only under the terms
at https://www.gnu.org/licenses/translations.html.

この参考訳は、変更の有無を問わず、
https://www.gnu.org/licenses/translations.html に掲げる条件に従う場合に限り公表できる。


\begin{center}
{\bf\large 前文（Preamble）}
\end{center}

本ライセンスの目的は、マニュアル、教科書その他の機能的かつ有用な文書を、自由という
意味で「フリー」にすること、すなわち、その文書を変更してもしなくても、また営利目的で
あっても非営利目的であっても、誰もが実際に複製し再頒布できる自由を確保することである。
第二に、本ライセンスは、他者が行った変更について責任を負うことなく、著者と出版者が
その仕事について適切な評価を受けられる手段を保全する。

本ライセンスは「コピーレフト」の一種であり、文書の派生著作物も同じ意味で自由でなければ
ならないことを意味する。これは、自由ソフトウェアのために設計されたコピーレフト・
ライセンスである GNU 一般公衆ライセンスを補完するものである。

本ライセンスは、自由ソフトウェアのマニュアルに利用するために設計された。自由ソフトウェア
には自由な文書が必要であり、自由なプログラムには、そのソフトウェアと同じ自由を提供する
マニュアルが伴うべきだからである。しかし、本ライセンスはソフトウェアのマニュアルに限られ
ない。主題が何であるか、また印刷書籍として出版されるかどうかを問わず、あらゆる文章著作物
に利用できる。とりわけ、教育または参照を目的とする著作物に本ライセンスを利用することを
推奨する。


\section{適用範囲と定義（APPLICABILITY AND DEFINITIONS）}
\label{fdl-section-applicability-and-definitions}

本ライセンスは、いかなる媒体によるものであっても、著作権者が本ライセンスの条件に従って
頒布できる旨の告知を置いたすべてのマニュアルその他の著作物に適用される。その告知は、
ここに定める条件の下で当該著作物を利用するための、全世界に及び、使用料を要せず、期間の
制限がないライセンスを与える。以下、\textbf{「文書」（Document）}とは、そのような
マニュアルまたは著作物をいう。公衆の誰もが被許諾者であり、本文では\textbf{「あなた」}
と呼ぶ。著作権法上の許諾を必要とする方法で著作物を複製、改変または頒布したとき、あなたは
本ライセンスを受諾したものとする。

文書の\textbf{「改変版」（Modified Version）}とは、文書またはその一部を含む著作物を
いい、一字一句そのままに複製したものか、改変したものか、または他言語に翻訳したものかを
問わない。

\textbf{「補遺部分」（Secondary Section）}とは、文書の出版者または著者と文書全体の主題
（または関連事項）との関係だけを扱い、その全体的主題に直接含まれ得る内容を一切含まない、
文書中の題名を有する付録または前付け部分をいう。（したがって、文書の一部が数学の教科書
である場合、補遺部分で数学を説明してはならない。）その関係は、主題または関連事項との
歴史的なつながりであってもよく、それらに関する法律上、商業上、哲学上、倫理上または政治上
の立場であってもよい。

\textbf{「変更不可部分」（Invariant Sections）}とは、文書が本ライセンスの下で公開される
旨を述べる告知において、その題名が変更不可部分の題名として指定されている一定の補遺部分を
いう。ある部分が上記の補遺部分の定義に適合しない場合、その部分を変更不可部分に指定すること
はできない。文書には変更不可部分が一つもなくてもよい。文書が変更不可部分を一つも特定して
いなければ、変更不可部分は存在しない。

\textbf{「カバーテキスト」（Cover Texts）}とは、文書が本ライセンスの下で公開される旨を
述べる告知において、表カバーテキストまたは裏カバーテキストとして列挙される短い文章をいう。
表カバーテキストは5語以内、裏カバーテキストは25語以内でなければならない。

文書の\textbf{「透過的」（Transparent）}複製物とは、仕様が一般公衆に利用可能な形式で
表現された機械可読な複製物であり、一般的なテキストエディタ、（画素からなる画像については）
一般的なペイントプログラム、または（図面については）広く利用できる作図エディタを用いて
文書を直接かつ容易に改訂するのに適し、かつテキスト整形プログラムへの入力、またはその入力
に適する各種形式への自動変換に適するものをいう。それ以外は透過的なファイル形式であっても、
そのマークアップまたはマークアップの欠如が、読者による後の改変を妨害または抑制するように
仕組まれている複製物は透過的ではない。相当量の文章に用いられる画像形式は透過的ではない。
「透過的」でない複製物を\textbf{「非透過的」（Opaque）}という。

透過的な複製物に適する形式の例には、マークアップを含まないプレーン ASCII、Texinfo 入力
形式、LaTeX 入力形式、一般に利用可能な DTD を使用する SGML または XML、および人が改変
することを想定して設計された標準準拠の簡素な HTML、PostScript または PDF がある。
透過的な画像形式の例には PNG、XCF および JPG がある。非透過的な形式には、専有ワード
プロセッサでしか読み書きできない専有形式、DTD または処理ツールが一般に利用できない SGML
または XML、および一部のワードプロセッサが出力だけを目的として生成する HTML、PostScript
または PDF がある。

印刷書籍における\textbf{「題扉」（Title Page）}とは、題扉そのものに加えて、本ライセンスが
題扉への掲載を要求する資料を読みやすく収めるために必要な後続のページをいう。題扉自体がない
形式の著作物では、「題扉」とは、本文の開始に先立つ、著作物の題名が最も目立って現れる箇所の
近くにある文章をいう。

\textbf{「XYZ と題された」（Entitled XYZ）}部分とは、文書の題名を有する小部分であって、
その題名が正確に XYZ であるか、または XYZ を他言語に翻訳した文章の後の括弧内に XYZ を
含むものをいう。（ここで XYZ は、\textbf{「Acknowledgements」}、
\textbf{「Dedications」}、\textbf{「Endorsements」}または\textbf{「History」}など、
以下で言及する特定の部分名を表す。）文書を改変するとき、そのような部分の
\textbf{「題名を保存する」（Preserve the Title）}とは、この定義に従ってその部分が
「XYZ と題された」部分であり続けることをいう。

文書は、本ライセンスが文書に適用される旨の告知の近くに保証否認警告を含めてもよい。
それらの保証否認警告（Warranty Disclaimers）は、保証を否認する限りにおいて、
参照によって本ライセンスに組み込まれたものとみなされる。保証否認警告が持ち得るその他の
含意はすべて無効であり、本ライセンスの意味に何ら影響を及ぼさない。


\section{逐語的複製（VERBATIM COPYING）}
\label{fdl-section-verbatim-copying}

あなたは、営利目的か非営利目的かを問わず、いかなる媒体でも文書を複製し頒布できる。ただし、
すべての複製物に本ライセンス、著作権表示および本ライセンスが文書に適用される旨の許諾告知を
再現し、本ライセンスの条件に他のいかなる条件も追加してはならない。作成または頒布する複製物
の閲覧またはさらなる複製を妨げ、もしくは制御するために技術的手段を用いてはならない。ただし、
複製物と引き換えに対価を受け取ることはできる。十分に多数の複製物を頒布する場合は、第3項の
条件にも従わなければならない。

上記と同じ条件の下で、複製物を貸し出し、公に展示することもできる。


\section{大量の複製（COPYING IN QUANTITY）}
\label{fdl-section-copying-in-quantity}

文書の印刷された複製物（または通常は印刷された表紙を備える媒体による複製物）を100部を
超えて出版し、文書の許諾告知がカバーテキストを要求する場合、複製物を、これらすべての
カバーテキストを明瞭かつ読みやすく掲載した表紙で包まなければならない。すなわち、表表紙には
表カバーテキストを、裏表紙には裏カバーテキストを掲載する。両方の表紙は、あなたがこれらの
複製物の出版者であることも明瞭かつ読みやすく示さなければならない。表表紙には完全な題名を、
題名中のすべての語が同じように目立ち、見えるように掲げなければならない。それに加えて、
表紙に他の資料を加えてもよい。変更が表紙だけに限られ、文書の題名を保存し、これらの条件を
満たす複製は、その他の点では逐語的複製として扱うことができる。

いずれかの表紙に必要な文章が長すぎて読みやすく収まらない場合、列挙されたものの先頭から
（合理的に収まるだけ）実際の表紙に置き、残りを隣接するページに続けるべきである。

文書の非透過的複製物を100部を超えて出版または頒布する場合は、機械可読な透過的複製物を
各非透過的複製物と共に添付するか、または各非透過的複製物の中もしくはそれと共に、一般の
ネットワーク利用者が、公に標準的なネットワークプロトコルを用いて、追加物のない完全な
透過的複製物をダウンロードできるコンピュータネットワーク上の場所を明示しなければならない。
後者を選択する場合、その版の非透過的複製物を公衆に最後に頒布した時（直接頒布した場合も、
代理人または小売業者を通して頒布した場合も含む）から少なくとも1年後まで、その透過的複製物
が明示した場所で引き続き利用可能であることを確保するため、非透過的複製物の大量頒布を開始
するときに、合理的で慎重な措置を講じなければならない。

多数の複製物を再頒布する十分前に、文書の著者が文書の更新版を提供できるよう著者に連絡する
ことが望まれるが、これは義務ではない。


\section{改変（MODIFICATIONS）}
\label{fdl-section-modifications}

あなたは、上記第2項および第3項の条件の下で文書の改変版を複製し頒布できる。ただし、改変版を
正確に本ライセンスの下で公開し、改変版に文書の役割を担わせることにより、その複製物を所有
するすべての者に改変版の頒布と改変を許諾しなければならない。さらに、改変版について次の
すべてを行わなければならない。

\begin{itemize}
\item[A.]
   題扉（および表紙があればその表紙）に、文書の題名および以前の版の題名とは異なる題名を
   使用すること。以前の版があれば、その題名は文書の「History」部分に列挙されているはずで
   ある。以前の版の原出版者が許可した場合は、その版と同じ題名を使用してもよい。

\item[B.]
   改変版の改変部分の著作に責任を負う一人以上の個人または団体を、著者として題扉に列挙し、
   文書の主要著者を少なくとも五人（主要著者が五人未満であれば全員）併記すること。ただし、
   それらの著者がこの要件を免除した場合を除く。

\item[C.]
   改変版の出版者の名称を、出版者として題扉に記載すること。

\item[D.]
   文書のすべての著作権表示を保存すること。

\item[E.]
   他の著作権表示の近くに、あなたが行った改変に対する適切な著作権表示を追加すること。

\item[F.]
   著作権表示の直後に、下記の付記に示す形式で、本ライセンスの条件の下で改変版を利用する
   ことを公衆に許可する許諾告知を含めること。

\item[G.]
   その許諾告知において、文書の許諾告知に掲げられた変更不可部分と必要なカバーテキストの
   完全な一覧を保存すること。

\item[H.]
   変更されていない本ライセンスの複製物を含めること。

\item[I.]
   「History」と題された部分を保存し、その題名を保存し、題扉に記載された改変版の題名、
   年、新たな著者および出版者を少なくとも記載した項目をそこに追加すること。文書に
   「History」と題された部分がなければ、題扉に記載された文書の題名、年、著者および出版者
   を記載する部分を作成し、その後に、前文に従って改変版を説明する項目を追加すること。

\item[J.]
   文書の透過的複製物を公衆が利用できるよう文書中に指定されたネットワーク上の場所があれば
   それを保存し、同様に、文書が基にした以前の版について文書中に指定されたネットワーク上の
   場所も保存すること。これらは「History」部分に置いてもよい。その著作物が文書自体より
   少なくとも4年前に出版された場合、またはその場所が指す版の原出版者が許可した場合は、
   ネットワーク上の場所を省略してもよい。

\item[K.]
   「Acknowledgements」または「Dedications」と題された部分について、その部分の題名を
   保存し、そこに記された各貢献者への謝辞および／または献辞の実質と語調をすべて保存すること。

\item[L.]
   文書のすべての変更不可部分を、その文章と題名を変更せずに保存すること。部分番号または
   それに相当するものは、部分の題名の一部とはみなされない。

\item[M.]
   「Endorsements」と題された部分をすべて削除すること。そのような部分を改変版に含めては
   ならない。

\item[N.]
   既存の部分を「Endorsements」と題されるよう改題したり、その題名を変更不可部分の題名と
   競合するものに改題したりしないこと。

\item[O.]
   すべての保証否認警告を保存すること。
\end{itemize}

改変版に、補遺部分に該当し、文書から複製した資料を含まない新たな前付け部分または付録が
含まれる場合、あなたは任意に、その一部または全部を変更不可部分に指定できる。そのためには、
それらの題名を改変版の許諾告知にある変更不可部分の一覧に追加する。それらの題名は、他の
すべての部分の題名と異なっていなければならない。

改変版について諸者が行う推薦だけを含む場合、「Endorsements」と題された部分を追加できる。
たとえば、査読についての声明や、ある組織が本文を標準の権威ある定義として承認した旨の声明
などである。

改変版のカバーテキスト一覧の末尾に、5語以内の一文を表カバーテキストとして、25語以内の
一文を裏カバーテキストとして追加できる。一つの主体が（または一つの主体の手配により）追加
できるのは、表カバーテキスト一文と裏カバーテキスト一文だけである。文書の同じ表紙に、以前
あなたが追加した、またはあなたが代理する同じ主体の手配により追加されたカバーテキストが
すでに含まれている場合、別のものを追加してはならない。ただし、以前のカバーテキストを追加
した出版者から明示的な許可を得れば、以前のものを置き換えることができる。

文書の著者および出版者は、本ライセンスによって、改変版の宣伝に自己の名称を使用することも、
改変版を推薦していると主張または示唆することも許可していない。


\section{文書の結合（COMBINING DOCUMENTS）}
\label{fdl-section-combining-documents}

あなたは、上記第4項で改変版について定める条件に従い、文書と本ライセンスの下で公開された
他の文書とを結合できる。ただし、すべての原文書の変更不可部分を改変せずに結合作品に含め、
そのすべてを結合作品の許諾告知に変更不可部分として列挙し、かつ、それらのすべての保証否認
警告を保存しなければならない。

結合作品には本ライセンスの複製物を一つだけ含めればよく、同一の変更不可部分が複数ある場合、
一つの複製物で置き換えてよい。同じ題名で内容が異なる変更不可部分が複数ある場合、各部分の
題名の末尾に、判明していればその部分の原著者または出版者の名称を、そうでなければ一意な番号
を括弧内に加えて、各題名を一意にしなければならない。結合作品の許諾告知にある変更不可部分の
一覧中の題名にも、同じ調整を行わなければならない。

結合に際しては、各原文書中の「History」と題された部分を結合して、一つの「History」と
題された部分を作らなければならない。同様に、「Acknowledgements」と題された部分同士、
「Dedications」と題された部分同士を結合しなければならない。「Endorsements」と題された
部分はすべて削除しなければならない。

\section{文書の収集（COLLECTIONS OF DOCUMENTS）}
\label{fdl-section-collections-of-documents}

文書と、本ライセンスの下で公開された他の文書とからなる収集物を作成し、各文書にある本ライ
センスの複製物を、収集物に含めた一つの複製物で置き換えることができる。ただし、その他の
すべての点において、各文書の逐語的複製に関する本ライセンスの規則に従わなければならない。

このような収集物から一つの文書を取り出し、本ライセンスの下で単独に頒布することができる。
ただし、取り出した文書に本ライセンスの複製物を挿入し、その文書の逐語的複製に関するその他
すべての点において本ライセンスに従わなければならない。


\section{独立した著作物との集積（AGGREGATION WITH INDEPENDENT WORKS）}
\label{fdl-section-aggregation-with-independent-works}

文書またはその派生著作物と、他の別個かつ独立した文書または著作物とを、記憶媒体または頒布
媒体の一巻の中または上にまとめた編集物は、その編集により生じる著作権が、個々の著作物が許可
する範囲を超えて編集物の利用者の法的権利を制限するために用いられない場合、
「集積物」（aggregate）と呼ばれる。文書が集積物に含まれるとき、本ライセンスは、
それ自体が文書の派生著作物でない集積物中の他の著作物には適用されない。

第3項のカバーテキスト要件が文書のこれらの複製物に適用され、かつ文書が集積物全体の半分未満
である場合、文書のカバーテキストは、集積物の中で文書を囲む表紙、または文書が電子的形式で
ある場合にはそれと同等の電子的表紙に置いてよい。それ以外の場合、集積物全体を囲む印刷された
表紙に掲載しなければならない。


\section{翻訳（TRANSLATION）}
\label{fdl-section-translation}

翻訳は改変の一種とみなされるため、第4項の条件の下で文書の翻訳を頒布できる。変更不可部分を
翻訳で置き換えるには、その著作権者の特別な許可が必要である。ただし、変更不可部分の原版に
加えて、その一部または全部の翻訳を含めることができる。本ライセンス、文書中のすべての許諾
告知およびすべての保証否認警告について、その翻訳を含めることができる。ただし、本ライセンス
の英語原版ならびにそれらの告知および否認警告の原版も含めなければならない。本ライセンス、
告知または否認警告の翻訳と原版との間に不一致がある場合は、原版が優先する。

文書に「Acknowledgements」、「Dedications」または「History」と題された部分がある場合、
その題名を保存するという第4項の要件（第1項に定義）は、通常、実際の題名を変更することを
必要とする。


\section{終了（TERMINATION）}
\label{fdl-section-termination}

本ライセンスに明示的に定める場合を除き、文書を複製、改変、サブライセンスまたは頒布しては
ならない。それ以外に文書を複製、改変、サブライセンスまたは頒布しようとする試みはすべて
無効であり、本ライセンスに基づくあなたの権利を自動的に終了させる。ただし、あなたから本
ライセンスに基づく複製物または権利を受け取った者のライセンスは、その者が完全な遵守を続ける
限り、終了しない。


\section{本ライセンスの将来の改訂（FUTURE REVISIONS OF THIS LICENSE）}
\label{fdl-section-future-revisions-of-this-license}

Free Software Foundation は、GNU 自由文書ライセンスの新たな改訂版を随時公表することが
ある。それらの新版は現在の版と理念を同じくするが、新たな問題または懸念に対処するため細部が
異なることがある。https://www.gnu.org/copyleft/ を参照せよ。

本ライセンスの各版には、他の版と区別するための版番号が付される。文書が、本ライセンスの特定
の版番号「またはそれ以降の版」が適用されると指定している場合、あなたは、指定された版または
Free Software Foundation が（草案としてではなく）公表したそれ以降の版のいずれかの条件に
従うことを選択できる。文書が本ライセンスの版番号を指定していない場合、Free Software
Foundation がこれまでに（草案としてではなく）公表したいずれの版も選択できる。


\section{付記：本ライセンスをあなたの文書に利用する方法（ADDENDUM: How to use this License for your documents）}
\label{fdl-section-addendum-how-to-use-this-license-for-your-documents}

あなたが作成した文書に本ライセンスを利用するには、本ライセンスの複製物を文書に含め、題扉
の直後に次の著作権表示と許諾告知を置く。

\bigskip
\begin{quote}
    Copyright \copyright  YEAR  YOUR NAME.
    Permission is granted to copy, distribute and/or modify this document
    under the terms of the GNU Free Documentation License, Version 1.2
    or any later version published by the Free Software Foundation;
    with no Invariant Sections, no Front-Cover Texts, and no Back-Cover Texts.
    A copy of the license is included in the section entitled "GNU
    Free Documentation License".
\end{quote}
\bigskip

変更不可部分、表カバーテキストおよび裏カバーテキストがある場合は、
「with...Texts.」の行を次のように置き換える。

\bigskip
\begin{quote}
    with the Invariant Sections being LIST THEIR TITLES, with the
    Front-Cover Texts being LIST, and with the Back-Cover Texts being LIST.
\end{quote}
\bigskip

変更不可部分はあるがカバーテキストはない場合など、三者の別の組合せについては、状況に
合うように上記二つの選択肢を組み合わせる。

文書に重要なプログラムコードの例が含まれる場合、自由ソフトウェアで利用できるように、
GNU 一般公衆ライセンスなど、選択した自由ソフトウェア・ライセンスの下でも、それらの例を
並行して公開することを推奨する。

%---------------------------------------------------------------------
